\section{Numerical tests}
\label{sec:tests}

\begin{figure}
\includegraphics[width=0.78\textwidth]{images/stout_test_B57}
\caption[figA]{The effective energy $aE_{\rm eff}(r)$ defined in 
 Eq.~(\ref{eq:Eeff}) for a static quark-antiquark pair separated by a 
 distance $r=5a$ for several levels of smearing $n_\rho=1,5,$ and $20$
 against the smearing parameter $\rho$.
 These results were obtained on a $12^4$ lattice using the Wilson gauge
 action with coupling $\beta=5.7$. 
 Curves labeled $A_{n_\rho}$ indicate results using the spatially-isotropic
 three-dimensional version of the analytic stout link
 smearing scheme with $n_\rho=1,5,$ and $20$ levels, while the curves
 labeled $B_{n_\rho}$ show the results for links smeared using
 Eq.~(\ref{eq:Ufuzz}) with the projection method of Ref.~\cite{su3projectA}.
\label{fig:stout_veff}}
\end{figure}

The efficacy of stout links both as a gluonic-operator smearing algorithm and
for reducing the effects of ultraviolet gluon modes in short-distance quantities
was tested in several Monte Carlo simulations.

The energy of gluons in the presence of a static quark-antiquark pair
separated by a distance $r$ can be extracted from $r\times t$ Wilson
loops $W(r,t)$ as the temporal extent $t$ becomes large.  The Wilson
loop can be viewed as the correlation function of a gauge-invariant
operator consisting of a static quark-antiquark pair connected
by a product of link variables following a straight-line path between the
quark and the antiquark.  Extraction of the lowest energy can be done much
more reliably if the couplings of the quark-antiquark-gluon operator
with higher-lying states are small.  This can be achieved by utilizing
a product of smeared links connecting the quark and the antiquark, instead
of the original link variables.  In other words, the ground state energy of
gluons in the presence of a static quark-antiquark pair can be more
reliably determined from Wilson loops $\tilde{W}(r,t)$ constructed using
smeared spatial links in the $r$-link paths on the initial and final
time slices.  A measure of how well the couplings to the higher levels are
reduced is the effective energy for a single time step, defined by
\begin{equation}
  a E_{\rm eff}(r) = -\ln\left(\frac{\tilde{W}(r,a)}{\tilde{W}(r,0)}
 \right).
 \label{eq:Eeff}
\end{equation}
Under unexceptional circumstances for a lattice gauge action with a
positive-definite transfer matrix, such as the Wilson action,
reducing the excited-state contamination tends to lower the
effective energy given above.


\begin{figure}
\includegraphics[width=0.78\textwidth]{images/stout_test_B62}
\caption[figB]{The effective energy $aE_{\rm eff}(r)$ defined in 
 Eq.~(\ref{eq:Eeff}) for a static quark-antiquark pair separated by a 
 distance $r=10a$ for several levels of smearing $n_\rho=1,5,$ and $20$
 against the smearing parameter $\rho$.
 Results were obtained on a $24^4$ lattice using the Wilson gauge action
 with coupling $\beta=6.2$. 
 Curves labeled $A_{n_\rho}$ indicate results using the spatially-isotropic
 three-dimensional version of the analytic stout link
 smearing scheme with $n_\rho=1,5,$ and $20$ levels, while the curves
 labeled $B_{n_\rho}$ show the results for links smeared using
 Eq.~(\ref{eq:Ufuzz}) with the projection method of Ref.~\cite{su3projectA}.
\label{fig:stout_veff2}}
\end{figure}

The effective energy defined in Eq.~(\ref{eq:Eeff}) was used to test
the ability of the stout link smearing scheme to reduce excited-state
contamination in gluonic operators.   Results for $r=5a$ on a $12^4$
lattice using the Wilson gauge action with coupling $\beta=5.7$ are shown
in Fig.~\ref{fig:stout_veff}, and results for $r=10a$ on a larger $24^4$
lattice with $\beta=6.2$ are shown in Fig.~\ref{fig:stout_veff2}.
The results are compared to those obtained with links smeared using
Eq.~(\ref{eq:Ufuzz}) and the projection method of Ref.~\cite{su3projectA}.

Both figures show that link smearing can dramatically reduce
contamination from the high-lying modes of the theory.  Consider first
the results obtained using the spatially-isotropic three-dimensional
version of the stout link smearing scheme $(\rho_{jk}=\rho,
\ \rho_{4\mu}=\rho_{\mu 4}=0)$.  For a given
number of smearing iterations $n_\rho$, the effective energy decreases
initially as the smearing parameter $\rho$ is increased from zero.  
Eventually an optimal value for the smearing
parameter is reached at which the effective energy is minimized.
This optimal value decreases as the number of smearing levels $n_\rho$ is
increased, and the reduction in the effective energy at this optimal
$\rho$ value is substantial.  Further increasing $\rho$ beyond its
optimal value then results in a sharply increasing effective energy.
The rapidity of both the fall and rise of the effective energy
about its minimum is more pronounced for larger $n_\rho$. Also, the
minimum value decreases as $n_\rho$ increases until eventually a
saturation point is reached beyond which no additional reduction occurs.

The results obtained with links smeared using Eq.~(\ref{eq:Ufuzz}) and
the projection method of Ref.~\cite{su3projectA} display essentially
the same trends, except that each minimum in the effective energy is much
broader.  The increased sensitivity of the stout link smearing scheme
to the parameter $\rho$ is not surprising since $\rho$ occurs inside
an exponential function.   However, it is important to note that both
smearing methods produce nearly the same minimum values of the effective
energy.  The stout link smearing scheme is just as effective at reducing
excited-state contamination in gluonic operators as the standard
link fuzzing scheme in current use, although more careful tuning is
necessary due to the increased sensitivity to the smearing parameter $\rho$.
Note that this analytic link smearing method has already been successfully
applied in computing the spectrum of torelon excitations\cite{torelons}.

\begin{figure}
\includegraphics[width=0.78\textwidth]{images/stout4D_plaq_B57}
\caption[figC]{The mean smeared plaquette against the smearing parameter
$\rho$. These results were obtained using the Wilson gauge action on a
$12^4$ lattice with coupling $\beta=5.7$. Curves labeled $A_{n_\rho}$
indicate results obtained using the isotropic four-dimensional
version of the analytic stout link
smearing scheme with $n_\rho=1$ and $5$ levels, while the curves
labeled $B_{n_\rho}$ show the results with links smeared using
Eq.~(\ref{eq:Ufuzz}) and the projection method of Ref.~\cite{su3projectA}.
\label{fig:stout_plaq}}
\end{figure}

The main purpose of using smeared links in lattice gauge and fermion
actions is to reduce discretization effects caused by the ultraviolet
modes in the lattice theory.  Often such effects are dominated by
large contributions from tadpole diagrams.  A simple measure of
how well a link smearing scheme can reduce artifacts from tadpole
diagrams is the mean smeared plaquette.  A value of the mean smeared
plaquette near unity indicates a substantial reduction of the tadpole
contributions.

Results for the mean smeared plaquette on a $12^4$ lattice using the
Wilson gauge action with coupling $\beta=5.7$ are shown in
Fig.~\ref{fig:stout_plaq}.  For isotropic four-dimensional versions
$(\rho_{\mu\nu}=\rho)$ of both smearing methods, the mean
smeared plaquette initially increases towards unity as $\rho$
increases from zero until a maximum is reached.  The two
smearing methods produce nearly the same maximum value.  As $\rho$ is
increased further, the mean smeared plaquette in the stout link scheme
quickly begins to fall, whereas the standard smearing scheme changes
little.  As $n_\rho$ initially increases from zero, the maximum value
of the mean plaquette also increases, but eventually a saturation point
is reached.  Note that the maximum values for $n_\rho=5$ are very near to
unity, suggesting a dramatic reduction of tadpole contributions. 
In summary, the analytic stout link smearing scheme is observed to be just
as efficacious as the standard smearing scheme in reducing discretization
effects from the ultraviolet gluon modes, but with an increased sensitivity
to the smearing parameter $\rho$.

